\begin{table}[htbp]
\centering
\caption{Structural Break Tests and Event-Window Analysis of Searcher Volume}
\label{tab:structural_breaks}
\begin{tabular}{llrl}
\toprule
Test / Methodology & Test Statistic & p-value / Percentile & Break Dates Detected / Notes \\
\midrule
Known-Date Chow Test: Dencun (2024-03-13) & F = 7.25 & p = 7.6150e-04 & Classical Chow (Descriptive; non-iid AR(1) errors) \\
Known-Date Chow Test: Pectra (2025-05-07) & F = 94.42 & p = 1.1102e-16 & Classical Chow (Descriptive; non-iid AR(1) errors) \\
Quandt-Andrews Unknown-Date Scan & sup-F = 212.94 & Bootstrap p = 0.0000 & Argmax date: 2025-07-10 (15% trimming, 14d block boot) \\
Ruptures Binary Segmentation (BIC Capped) & L2 Cost / BIC & --- & Descriptive segmentation on log volume (max 3 breaks): 2025-05-05 (index 490), 2025-07-09 (index 555), 2025-10-17 (index 655) \\
Event Window Ratio: Dencun (±60 Days) & Post/Pre Ratio = 1.21x & 77.5-th percentile & Compared to 200 random placebo dates (seed 42) \\
Event Window Ratio: Pectra (±60 Days) & Post/Pre Ratio = 1.37x & 90.0-th percentile & Compared to 200 random placebo dates (seed 42) \\
\bottomrule
\end{tabular}
\vspace{0.2cm}
\begin{minipage}{\linewidth}
\small
\textit{Notes:} This table summarizes structural break tests and event-window volume shifts around Ethereum upgrades on the bot-side daily series ($N=731$). \textbf{Non-iid Error Caveat:} Classical Chow tests assume independent and identically distributed errors, which daily volume violates ($AR(1) \approx 0.78$). Consequently, the known-date Chow tests are explicitly labeled as descriptive; primary formal inference is provided by the HAC-robust Wald F-tests on step and slope terms in Spec (2) of Table 6. The Quandt-Andrews scan reports the maximum Chow F over the central 70\% sample, with empirical p-value obtained via a 999-replication circular moving-block bootstrap (block length = 14 days, seed 42). Ruptures segmentation is conducted descriptively on log daily volume with binary segmentation capped at a maximum of 3 breaks to prevent over-segmentation on a trending series. Event window post/pre ratios report mean daily volume shifts in $\pm 60$-day windows, evaluated against a placebo distribution of 200 random non-event dates (seed 42); note that underlying trend growth confounds raw window ratios.
\end{minipage}
\end{table}
